\section{Unbounded optimized narrowness and completion of the proof}
\label{sec:unbounded-narrowness}

The remaining case of the second model is excluded by the following
proposition. Here \(\ell\) is the infimum, over maximizing sequences,
of the terminal narrowness supremum defined in
\Cref{sec:critical-paths}.

\begin{proposition}[Exclusion of unbounded optimized narrowness]
 \label{prop:unbounded-improvement}
 In the second model's extremal setup of \Cref{prop:critical-path},
 the global optimized narrowness \(\ell\) is finite.
\end{proposition}

Assume, for a contradiction, that \(\ell=\infty\).
We will construct terminal-known atlases at a fixed positive depth
whose time exponents tend to zero. Completing these atlases to the
terminal depth contradicts the positive critical time rate \(k\).
The construction has three stages. Two nested true systems first
give thin charts valid for the whole time interval. A further
normalization makes the logarithmic exponents of their mass-to-area
ratios tend to zero. These packets then satisfy
the scalar projection theorem: its quadratic fits can be combined
into an actual known atlas with the required depth gain and coverage.

Throughout this section, \(K_0=N_0^{o(1)}\) denotes a subpower factor
within a fixed exact problem; it may depend on that problem.
After taking the outer diagonal, \(K=N^{o(1)}\) denotes a subpower
factor along that diagonal. The latter controls the scalar analysis,
whereas the output atlas must retain the former type of bound.

\subsection{Nested true systems}

Take a true intermediate system on a maximizing critical path,
normalize through it as in \Cref{lem:normalization}, and rescale the
remaining depth interval to length one.
Denote the resulting maximizing sequence by \(E_i\). In this section,
\emph{original coordinates} means these coordinates of \(E_i\), before
the further localizations below. The own-depth knowledge and
saturation of the initial true system give
\[
 n_i(0)\sim_{\log}1.
\]
Consequently the relative-profile conclusion of
\Cref{prop:critical-path} gives
\begin{equation}
 \label{a08:initial-profiles}
 f_i(r)\longrightarrow dr
 \quad\hbox{uniformly for }0\le r\le1,\qquad
 d_i\longrightarrow d=d_*<1,\qquad
 H(E_i)\longrightarrow k>0.
\end{equation}
Here \(n_i(r)=N_0^{-f_i(r)}\) in logarithmic order, and
\(E_i\in C(d_i)\). The fixed angular parameters are denoted by
\(S,\eta>0\).

The original trajectories satisfy \(|y|+|V|\le K_0\), and their native
signal \(X=P_*(t,y,w)\) has hidden derivative scale \(B\).  At an exact
base \((t,y)\), the conditions
\[
 |P_*|\le K_0,\qquad |(P_*)_w|\ge B/K_0
\]
on incidences require only \(K_0\) hidden bins of width \(1/B\).
Indeed a quadratic has boundedly many monotonicity intervals, and on
the tested part its inverse image of an interval of length \(O(K_0)\)
has total length \(O(K_0^2/B)\).  Together with the zero-depth density
bound in \Cref{a08:initial-profiles}, this gives a base-density ceiling
\(K_0\) in each world after removal of atypical entries.

Choose parameters in the order
\begin{equation}
 \label{a08:depth-choice}
 0<u<\eta,\qquad 0<s_1<s_2<1,\qquad
 (d+2k)s_2<\frac{Su}{4}.
\end{equation}
The segment construction in \Cref{prop:critical-path} supplies nested
true nodes \(Q_1,Q_2\), with depths \(r_{j,i}\to s_j\) and common
homogenized time lengths \(q_j\), such that
\begin{equation}
 \label{a08:two-node-scales}
 q_j=N_0^{-h_{j,i}+o(1)},\quad h_{j,i}\to ks_j,\qquad
 g_j=N_0^{-\alpha_{j,i}},\qquad
 \alpha_{2,i}-\alpha_{1,i}\longrightarrow\infty.
\end{equation}
Here \(q_j^2g_j\) is the spatial determinant in original coordinates,
up to actual subpower factors.  To obtain the last limit, first normalize
at the earlier node and choose the later endpoint as in the critical-path
construction.  The intervening problem, divided by its actual depth
length, is again maximizing.  The global hypothesis \(\ell=\infty\)
therefore gives, by part~4 of \Cref{prop:critical-path}, true systems
whose horizon tends to \(k\) and whose narrowness tends to infinity.
Multiplication by the intervening depth length, which tends
to \(s_2-s_1>0\), preserves divergence when these systems are lifted.

All labels include their required path histories.  We retain their
full-layer trajectory assignments, so that, within each original world
\(\gamma\),
\begin{equation}
 \label{a08:old-layer-masses}
 \nu_\gamma(Q_j)\sim_{\log}
 n_i(r_{j,i})q_j^2g_j.
\end{equation}
The subsequent path incidences are subsets of these assignments.
Assignments at one node and in one time block are disjoint; an earlier
floor is not asserted for every later restriction.

Both labels are known by \(p_L=p_1\).  The angular fullness statement
therefore makes their largest spatial widths \(q_j\) up to actual
subpowers.  Write \(D_j\) for the spatial matrix, \(y_j(t)\) for its
affine center, and \(n_j\) for a unit left minor axis.  The position
bounds and their endpoint consequences give
\begin{equation}
 \label{a08:normal-strips}
 \begin{aligned}
 |n_j\cdot(y-y_j(t))|&\le K_0q_jg_j
       &&\text{throughout the \(Q_j\)-block},\\
 |n_j\cdot(V-y_j')|&\le K_0g_j .
 \end{aligned}
\end{equation}
The full center velocities are bounded by \(K_0\), since the trajectories
have bounded velocities and fit the chart throughout a nondegenerate
block.

At a common path incidence,
\begin{equation}
 \label{a08:nested-angle}
 |\sin\angle(n_1,n_2)|\le K_0g_1 .
\end{equation}
In fact \(D_2=D_1R\), where the relative matrix satisfies
\(\|R\|\le K_0q_2/q_1\).  Thus
\(\|n_1^{\mathsf T}D_2\|\le K_0q_2g_1\).
The major axis of \(D_2\) has length at least \(q_2/K_0\), which proves
\Cref{a08:nested-angle}.

\subsection{Localization for the whole time interval}

The true charts control each trajectory only on their own blocks.
We first compare two occurrences of the later chart on the same
trajectory.  Their nearly parallel normal strips will determine one
thin chart valid for the entire original time interval.

\begin{proposition}[Full-time localization]
 \label{prop:full-time-localization}
 In the setting of \Cref{a08:depth-choice,a08:two-node-scales}, a permitted
 further selection admits full-time charts \(C\) with tangential width
 one and normal width
 \begin{equation}
  \label{a08:C-width}
  g_C=g_2N_0^{2u}.
 \end{equation}
 Their labels are known by the original \(p_1\) with actual subpower
 lists.  The native test remains \(P_*\).  If \(n_C\) and \(y_C(t)\)
 are the normal and affine center of \(C\), then every retained
 \(Q_1\)-incidence satisfies, after a consistent choice of signs,
 \begin{equation}
  \label{a08:C-Q1-comparison}
  |n_1-n_C|\le K_0g_1,\qquad
  |n_C\cdot(y_1'-y_C')|\le K_0g_1 .
 \end{equation}
 All assignments, geometric tests and coordinate changes have tempered
 bounds within each fixed \(E_i\).
\end{proposition}

\begin{proof}
Sample two independent times \(t_0,t\) on the same trajectory, using the
original trajectory and world priors, and require a current path
incidence at both times.  If the path has mass \(p\), the pair mass is
at least \(p^2\) by Jensen's inequality, and hence at least \(1/K_0\).

We first prove that, outside negligible pair mass,
\begin{equation}
 \label{a08:full-time-angle}
 |\sin\angle(n_2(t_0),n_2(t))|\le g_2N_0^u.
\end{equation}
Before counting labels, remove from the \(t\)-incidences the atypical
sampled entries for the terminal joint angular cap at depth \(u/2\),
the pure terminal lower density, and the zero-depth upper density.
Their removal costs negligible pair mass: the additional time has
length at most one, and the density tolerances can be relaxed slowly
relative to the subpower path mass.

For fixed \(t_0,\gamma\), saturation and disjointness give at most
\begin{equation}
 \label{a08:first-label-count}
 \frac{K_0}{n_i(r_{2,i})q_2^2g_2}
\end{equation}
possible first labels \(Q_2(t_0)\).  Fix one.  If
\Cref{a08:full-time-angle} fails, the two velocity strips in
\Cref{a08:normal-strips} intersect in a ball of radius at most
\(K_0N_0^{-u}\).  At each \(p_1(t)\), its actual list contains only
\(K_0\) possible second labels.  Thus only \(K_0\) angular bins of
side comparable to \(N_0^{-u/2}\) must be considered.  On the retained
entries their total numerator is at most
\(K_0N_0^{-Su/2}\) times the original pure \(p_1\) density.

To integrate this denominator, sum only over its permitted good
zero-depth ancestors.  At each exact base only \(K_0\) such hidden
bins occur, and their total density is at most \(K_0\).
Moreover, a trajectory assigned to the fixed first label stays, over
the whole unit interval, in its extrapolated transverse strip of width
\(K_0g_2\).  Its tangential coordinate remains bounded.  The area
available at time \(t\) is therefore at most \(K_0g_2\).
Multiplying by \Cref{a08:first-label-count}, and then integrating
the times and weighting the worlds, bounds the failure mass by
\begin{equation}
 \label{a08:angle-failure}
 \frac{K_0N_0^{-Su/2}}{q_2^2n_i(r_{2,i})}
 =
 N_0^{-Su/2+(d+2k)s_2+o_i(1)+o_{N_0}(1)}.
\end{equation}
This is power-small by \Cref{a08:depth-choice}.  The exceptional
density entries were removed before the sum, so their mass is not
multiplied by the number of first labels.

Fix \(t_0\) by averaging, retaining subpower mass of paired successes.
The actual retained selection may use the original tempered path and
the angle test; the analytical density masks are not required as output
predicates.  Assign each successful trajectory at \(t_0\) according to
the following binned data of \(Q_2(t_0)\): its unit normal, and the
intercept and velocity coefficients of its affine normal center.
Use bins of width \(g_C\), treating simultaneous sign reversal by
finitely many fixed orientation charts.  These bins, without the
identity of \(Q_2(t_0)\), are the \(C\)-label.

For each bin choose a representative unit normal and a representative
affine normal motion.  The transverse bound in
\Cref{a08:normal-strips} extrapolates to \(K_0g_2\) over unit time.
Rounding the normal and the two coefficients adds at most \(K_0g_C\),
because positions and velocities are bounded by \(K_0\).
A tangential column of length one and a normal column of length \(g_C\)
therefore give the required full-time geometry.  The native \(P_*\)
already has the full-block bounds for a depth-zero chart.

We verify terminal knowledge of the whole label, including its center
motion.  A later \(Q_2(t)\) in the list at \(p_1(t)\) has normal within
\(K_0g_2N_0^u=o(g_C)\) of the first normal, modulo sign.
Both labels constrain the same bounded affine trajectory.  Their
normal intercepts and normal velocities consequently differ by at most
\(K_0g_2N_0^u\) as well: compare each coefficient to the corresponding
coefficient of that trajectory, using
\Cref{a08:normal-strips}, and then account for the change of normal.
Each later label thus determines boundedly many neighboring bins for
all the \(C\)-data.  Composing with the actual \(Q_2(t)\)-list proves
that \(C\) is known by \(p_1(t)\) with actual subpower lists.

Since
\[
 \frac{g_C}{g_1}
   =N_0^{-(\alpha_{2,i}-\alpha_{1,i})+2u},
\]
the binning error is much smaller than \(g_1\) for large \(i\).
Combining with \Cref{a08:nested-angle} gives the first inequality in
\Cref{a08:C-Q1-comparison}.  At a retained occurrence, compare both
center velocities with the actual \(V\).  The \(Q_1\) normal error is
\(K_0g_1\), the \(C\) normal error is \(K_0g_C\), and changing the
normal costs at most \(K_0g_1\) because \(V-y_1'\) is bounded.
This proves the second inequality.  The relevant normals and center
velocities are constant parameters of the labels, so this comparison
holds throughout their blocks.

Slicing at a fixed real \(t_0\), the finite bin tests, the original path
predicates and the geometric angle test all have tempered complexity.
Although their bounds can grow with \(i\), they remain fixed polynomial
bounds along each exact sequence \(E_i\).
\end{proof}

\subsection{Balanced localized packets}

The full-time charts give a common thin direction, but their masses
need not yet be comparable to their spatial areas.  The next
normalization produces this comparison in the two-limit sense needed
for scalar projection.  Its time cost tends to zero, leaving room for
the fixed positive depth gain constructed below.

\begin{lemma}[Normalization and relative density cancellation]
 \label{a08:balanced-packets}
 After a further permitted selection and passage to subsequences, the
 charts \(C\) of \Cref{prop:full-time-localization} contain a system of
 charts \(A\), known by original \(p_1\), with common time length
 \[
 q_A=N_0^{-h_{A,i}+o(1)},\qquad h_{A,i}\longrightarrow0.
 \]
 Write \(D_A\) for their original spatial matrices and
 \(J_A=|\det D_A|\).  Their original assigned trajectory masses satisfy
 \begin{equation}
  \label{a08:balanced-density}
  \frac{\nu_\gamma(A)}{J_A}
   =N_0^{-db_i+\eps_i+o_{N_0}(1)}
   =N_0^{o_{i\to\infty}(1)+o_{N_0\to\infty}(1)},
  \qquad b_i\longrightarrow0,\quad \eps_i\longrightarrow0.
 \end{equation}
 The second equality is a two-limit statement.  It does not assert an
 actual subpower bound for fixed \(i\).

 If \(e_Z\) is a unit tangent to \(C\), then the systems may also be
 prepared so that
 \begin{equation}
  \label{a08:tangential-fullness}
  \|D_A^{\mathsf T}e_Z\|\ge q_A/K_0,\qquad
  \|D_A^{\mathsf T}n_C\|\le K_0g_Cq_A.
 \end{equation}
 Actual lists, native derivative bounds and geometric slacks remain
 subpower within each fixed exact problem.
\end{lemma}

\begin{proof}
Normalize spatially inside \(C\), keeping the time interval, hidden
coordinate, thickness and native \(P_*\) unchanged.
The depth-zero case of \Cref{lem:normalization} applies to this known
atlas without requiring saturation at depth zero.  Its new world
\((\gamma,C)\) has conditional trajectory prior
\(\nu_\gamma(\,\cdot\,|C)\).  Put
\[
 W=\sum_{\gamma,C}\omega_\gamma\nu_\gamma(C).
\]
Then \(K_0^{-1}\le W\le1\), and the normalized new world weights are
\(\omega_\gamma\nu_\gamma(C)/W\).
Multiplying the conditional law by its world weight restores the
original restricted trajectory law, divided only by \(W\).
The terminal density formula, within each world, is
\begin{equation}
 \label{a08:primed-terminal}
 n'_i(1)\sim_{\log}
 n_i(1)\frac{g_C}{\nu_\gamma(C)}.
\end{equation}
At intermediate depths only the one-sided comparison of
\Cref{lem:normalization} is needed.  To see the cap inheritance
explicitly, homogenize \(g_C/\nu_\gamma(C)=N_0^{t_i+o(1)}\).
Then \(f'_i(1)=f_i(1)-t_i\) and
\(f'_i(r)\ge f_i(r)-t_i\).  Hence
\[
 f'_i(1)-f'_i(r)\le f_i(1)-f_i(r)\le d_i(1-r).
\]
A primed velocity bin maps into an old velocity ball of its radius
times \(K_0\), since the norm of the \(C\)-matrix is bounded by
\(K_0\).  The old good terminal angular numerators therefore give
the same angular cap after the density change.  Thus the homogeneous
new problem \(E'_i\) belongs to \(C(d_i)\).

Every true terminal system in \(E'_i\) lifts through \(C\), so
\[
 H(E_i)\le H(E'_i)\le k(d_i).
\]
Since the outer sequence is maximizing, both endpoints tend to \(k\).
Consequently \(E'_i\) is also maximizing.  Its relative profiles have
the linear limits from \Cref{prop:critical-path}; no bound on
\(f'_i(0)\) or on another absolute primed exponent is needed.

Choose small positive depths \(b_i\to0\) slowly.  More explicitly, fix
\(b=1/m\), first take \(i\) late enough for the relative-profile
comparison at \(b,1\) with error at most \(1/m\), and choose true
packets in the length-\(b\) coarsening with horizon at most \(b+1/m\).
The universal initial bound \(H\le L\) provides these packets;
membership of the coarsening in the same cap class is not needed.
Let \(m=m(i)\) tend to infinity slowly enough to satisfy these
requirements together with the finite preparations already imposed.
Then
\begin{equation}
 \label{a08:primed-relative}
 f'_i(1)-f'_i(b_i)=d(1-b_i)+o_i(1),
\end{equation}
and the packet horizon tends to zero.  Saturate those primed true
packets and compose them with \(C\), including \(C\) in the resulting
label \(A\).

The primed spatial determinant of \(A\) is \(J_A/g_C\), up to actual
subpowers.  Saturation gives
\[
 \frac{\nu_\gamma(A)}{\nu_\gamma(C)}
   \sim_{\log}n'_i(b_i)\frac{J_A}{g_C}.
\]
Combining with \Cref{a08:primed-terminal} yields
\begin{equation}
 \label{a08:density-ratio}
 \frac{\nu_\gamma(A)}{J_A}
   \sim_{\log}
   n_i(1)\frac{n'_i(b_i)}{n'_i(1)}.
\end{equation}
The logarithmic exponent of the right-hand side is
\[
 -f_i(1)+f'_i(1)-f'_i(b_i)
   =-db_i+o_i(1),
\]
by \Cref{a08:initial-profiles,a08:primed-relative}.  This proves
\Cref{a08:balanced-density}.  In particular, an arbitrarily large
absolute primed density shift cancels in the quotient in
\Cref{a08:density-ratio}.

We record the mass bookkeeping for subsequent restrictions.
At the \(C\)-stage the factors \(\nu_\gamma(C)\) cancel conditional
normalization globally; only the actual subpower \(W^{-1}\) remains.
At the \(A\)-stage the relevant original budgets are
\begin{equation}
 \label{a08:A-budget}
 \sum_{\gamma,I,A}\omega_\gamma q_A\nu_\gamma(A)\le1.
\end{equation}
Here the sum includes the time blocks, and the inequality follows
from disjointness within each block.  Every later homogeneous
selection has subpower total success in its fixed exact problem.
An original exception of fixed power-small mass is therefore still
negligible relative to that success.  Worlds or labels with too small
a conditional success fraction may be discarded by summing
\Cref{a08:A-budget}.  We never require an exception bound uniformly
after conditioning on every individual tiny \(C\) or \(A\).

The label knowledge is also actual.  At \(p_1\), first use the actual
list for \(C\).  Given one such \(C\), its stored affine spatial map
determines the primed exact base.  The hidden coordinate and \(B\)
have not changed, so the old hidden bin of width \(N_0^{-1}/B\)
determines the coarser primed \(b_i\)-bin with at most two choices.
Compose this with the actual own-depth list of the primed true
packet.  If the two list sizes are \(L_C,L_{A'}\), the resulting
\(A\)-list has size at most \(2L_CL_{A'}=N_0^{o(1)}\) within fixed
\(i\), irrespective of the density shift in
\Cref{a08:primed-terminal}.

Composition gives the second inequality of
\Cref{a08:tangential-fullness} and also an original largest-axis upper
bound \(K_0q_A\).  Suppose on subpower retained mass that the first
inequality fails by a fixed power \(N_0^{-\eps}\).
The endpoint position tests then place the original velocity, for
each \(A\), in a ball whose radius is bounded by
\(K_0(N_0^{-\eps}+g_C)\).
Choose a fixed angular depth \(v>0\) smaller than \(\eps\), the allowed
angular range, and a positive power of smallness for \(g_C\).
The actual \(A\)-list at \(p_1\) requires only \(K_0\) corresponding
velocity bins.  On good terminal angular entries these bins carry
at most \(K_0N_0^{-Sv}\) times the original pure denominator.
Remove bad entries before summing and integrate over the old
observation.  The resulting power-small mass contradicts the
retained subpower mass.  Homogenization and slowly vanishing fixed
tolerances therefore give the first inequality of
\Cref{a08:tangential-fullness}.
\end{proof}

\subsection{The depth gain and the diagonal}

Retain the final tempered path in original worlds as the reference
incidence measure \(\mu_{\mathrm{ref}}\).  Write
\(\mu_{\mathrm{ref},\gamma}\) for its within-world measure, without
the world weight \(\omega_\gamma\), so that
\[
 \mu_{\mathrm{ref}}
   =\sum_\gamma\omega_\gamma\mu_{\mathrm{ref},\gamma},
 \qquad
 \|\mu_{\mathrm{ref}}\|=\mathfrak p_{\mathrm{path}}\ge K_0^{-1}.
\]
Here \(\|\mu\|\) denotes total mass.  The earlier full-layer
assignments and their prior masses remain fixed.
For later restricted measures, a subscript \(\gamma\) likewise
excludes the world weight.

We fix the depth gain and the losses allowed in coefficient matching
before selecting diagonal samples. Fix \(d<d_0<1\). Let
\(\eps_*\) be the tolerance in
\Cref{a02:local-graph-matching} for quadratic graphs on two base
variables, fixed comparison domains and dimension bound \(d_0\).
Choose fixed parameters
\[
 0<\theta<\min\{1,\eps_*/4\},\qquad
 2\theta<\eps<\eps_*.
\]
Let \(C_{\mathrm{match}}\ge1\) dominate the fixed exponents in
the list, coefficient, fit and marginal losses in the matching
application below.  These depend only on the fixed degrees and
comparison dimensions.  Choose
\[
 0<\delta<\min\left\{\frac{\theta s_1}{4},
                \frac{\eps s_1}{2C_{\mathrm{match}}}\right\},
\]
and put
\begin{equation}
 \label{a08:gain}
 c'=\frac{\theta s_1}{2}\in(0,1).
\end{equation}
For sufficiently large \(i\), an actual known atlas in \(E_i\) at
depth \(c'\) and time length \(q_A\) is impossible on an exact
subsequence.  Indeed \(h_{A,i}<kc'/2\), and normalization followed by a
true terminal system in the remaining problem would give
\begin{equation}
 \label{a08:forbidden-horizon}
 H(E_i)\le \frac{kc'}2+k(d_i)(1-c').
\end{equation}
The right-hand side tends to \(k-kc'/2<k\), contrary to
\Cref{a08:initial-profiles}.

Use the finite-sample exclusion argument of \Cref{a03:finite-guard} in the
following precise form.  Within fixed \(i\), fix a sufficiently large
polynomial complexity bound for the proposed outputs, and allow their
list sizes and geometric slacks to be bounded by a sufficiently large
fixed power of \(K_{\mathrm{path}}\log N_0\).  Here
\(K_{\mathrm{path}}\) collects actual subpower bounds for the old
problem and prepared path, including the reciprocal reference mass.
Exclude outputs covering at least
\begin{equation}
 \label{a08:required-output-mass}
 \mathfrak p_{\mathrm{path}}/\log N_0.
\end{equation}
If infinitely many samples in this fixed exact problem admitted such
outputs, their bounded complexity and actual subpower controls would
supply an exact known atlas on a subsequence, contradicting
\Cref{a08:forbidden-horizon}.  Thus one may sample \(N_0\) arbitrarily
late in this excluded range.

We take a diagonal \(N=N_0\to\infty\), \(i\to\infty\), sampling each
exact sequence late enough that the actual path factors and the
two-limit errors in \Cref{a08:balanced-density} are absorbed in
\(K=N^{o(1)}\).  The constants
\(u,s_1,s_2,\theta,\delta\) remain fixed.  All later finite geometric
operations have polynomial bounds depending only on the old
problem/path bounds and these fixed constants; the output complexity
bound above is chosen large enough for those operations before the
sample is taken.  The proof below supplies actual output lists and
slacks, rather than substituting this diagonal \(K\) for them.

Prepare analytical caps before sparse searches.  Remove atypical pure
densities at depths through \(s_1\), including \(0,r_{1,i}\), and bad
pure lower or joint angular upper entries at \(p_1\).
The loss is \(o(\mathfrak p_{\mathrm{path}})\).  One may use slowly
increasing finite depth grids, shrinking fixed exponent tolerances
and interpolation between neighboring widths.  At every fixed \(i\)
and tolerance the exceptional mass is power-small, so late sampling
makes all these preparations simultaneous.  In particular the pure
ceilings through \(s_1\) are \(KN^{-rd}\), and the ceiling at
\(r_{1,i}\) is \(Kn_i(r_{1,i})\).
These caps concern restricted measures and their specified good old
ancestors.  The analytical masks need not be encoded in the reference
path or in the final assignments.
Denote the resulting within-world submeasures by
\(\mu_{\mathrm{prep},\gamma}\le\mu_{\mathrm{ref},\gamma}\), and put
\[
 \mu_{\mathrm{prep}}
   =\sum_\gamma\omega_\gamma\mu_{\mathrm{prep},\gamma}.
\]
The next lemma makes the remaining light-label and light-entry
deletions and then fixes \(\mu_{\mathrm{prep}}\).  The projected laws
introduced afterward will all be restrictions of this one measure.

\subsection{Stable graphs on the earlier charts}

The scalar projection keeps time and the tangential coordinate
\(e_Z\cdot y\), together with \(X\). Its support count will use
the earlier \(Q_1\)-charts. We first give their native signal a finite
graph description with derivative bounds.  The lower masses needed here
are obtained before any conditioning on an \(A\)-chart or a projected
cell.

\begin{lemma}[Graph descriptions on the earlier charts]
\label{a08:earlier-graphs}
After analytical deletions of mass \(o(\mathfrak p_{\mathrm{path}})\),
each \(Q_1\)-label still meeting \(\mu_{\mathrm{prep}}\) has the
following properties.
\begin{enumerate}
\item It has a frozen earlier witness \(W_{\gamma,I,Q_1}\), supported
in its original time block and full spatial domain, with mass at least
\(q_1\nu_\gamma(Q_1)/K\).  These witnesses are disjoint within each
original block and satisfy the prepared zero-depth cap.
\item In its normalized base coordinates
\(\widehat u=((t-t_1)/q_1,D_1^{-1}(y-y_1(t)))\), where \(t_1\)
is the block's left endpoint, its retained incidences are assigned
to at most \(K\) pairs \((U,G)\).  Here \(U\) is a real box, called
the core, and \(G\) is a bounded-degree algebraic function,
holomorphic on a fixed-factor complex enlargement of \(U\).
Every assigned incidence has \(\widehat u\in U\) and satisfies
\[
 |X-G(\widehat u)|\le Ka,\qquad
 |G|+|\nabla G|\le K,
 \qquad a=N^{-s_1},
\]
where the value and derivative bounds for \(G\) hold on a
neighborhood of \(U\).
\end{enumerate}
The final prepared measure has mass
\(\|\mu_{\mathrm{prep}}\|=(1-o(1))\mathfrak p_{\mathrm{path}}\).
The first witnesses need not be submeasures of this final measure:
they are kept before the later core/graph deletions.
\end{lemma}

\begin{proof}
Delete \(Q_1\)-labels whose incidence mass under
\(\mu_{\mathrm{prep},\gamma}\), in their original block, is less than
\begin{equation}
 \label{a08:Q1-witness-floor}
 q_1\nu_\gamma(Q_1)/K.
\end{equation}
Choose the reciprocal-subpower threshold small enough relative to
the reference mass.  The loss is negligible because
\[
 \sum_{\gamma,I,Q_1}
   \omega_\gamma q_1\nu_\gamma(Q_1)\le1.
\]
Write \(W_{\gamma,I,Q_1}\) for each surviving restriction to its
\(Q_1\)-label and block at this stage.  Keep these measures as
earlier witnesses when subsequent restrictions decrease
\(\mu_{\mathrm{prep}}\).  The assigned prior
\(\nu_\gamma(Q_1)\) in
\Cref{a08:Q1-witness-floor} is the original full-layer prior in
\Cref{a08:old-layer-masses}.  The floor is global within \(Q_1\);
it is not a lower bound after conditioning on a particular \(A\).

We next describe the scalar signal by stable graphs in the normalized
coordinates of each surviving earlier chart.
Write \(z=Bw\), let \(P(\widehat u,z)\) be the true \(Q_1\)-test,
and set \(a_1^{\mathrm{fit}}=N^{-r_{1,i}}\sim_{\log,N}a\).
On the earlier retained incidences,
\begin{equation}
 \label{a08:local-test-bounds}
 |P|\le K_0a_1^{\mathrm{fit}},\quad
 K_0^{-1}\le |P_z|\le K_0,\quad
 |P_{zz}|\le K_0/a_1^{\mathrm{fit}}.
\end{equation}
The depth-\(r_{1,i}\) pure ceiling, the \(K_0\) compatible hidden
bins and the chart Jacobian give a base-density ceiling
\(Kq_1\nu_\gamma(Q_1)\) in \(\widehat u\)-coordinates.
The floor \Cref{a08:Q1-witness-floor} therefore implies a real base
projection of volume at least \(1/K\) in a box of size at most \(K_0\).

The discriminant
\[
 \mathcal D(\widehat u)=P_z^2-2P_{zz}P
\]
is independent of \(z\) and is a fixed-degree polynomial in
\(\widehat u\).  Its values on that projection are bounded by
\Cref{a08:local-test-bounds}.  Polynomial Remez and fixed-degree
coefficient bounds consequently bound \(\mathcal D\) and its first
derivatives by \(K\) on the needed complex enlargements of the
normalized box.

Choose a sufficiently large subpower cutoff \(K'\).
If \(|P_{zz}|\ge (a_1^{\mathrm{fit}}K')^{-1}\), the affine
quadratic center
\[
 z_{\mathrm c}(\widehat u)
   =-\frac{\text{coefficient of \(z\) in \(P\)}}{P_{zz}}
\]
approximates the observed \(z\) to \(Ka_1^{\mathrm{fit}}\), because
\(P_z=P_{zz}(z-z_{\mathrm c})\).
For the remaining tests choose \(K'\) large enough that
\(2|P_{zz}P|\ll |P_z|^2\) on the incidences.  Their discriminants
are positive and bounded below by \(1/K\) there.

Partition the normalized box into real core subboxes of
inverse-subpower radius.  Use fixed-factor complex enlargements
with an additional margin as comparison domains. For tests in the
small-curvature case, the derivative bound on \(\mathcal D\) permits
a radius such that its variation on the enlargement of every occupied
core is much smaller than its incidence lower bound. The discriminant
therefore never vanishes there, and the two simple roots are
holomorphic. For a linear equation the same construction keeps its
linear coefficient nonzero. In the large-curvature case use the affine
center \(z_{\mathrm c}\) on these same cores. At each occurrence the
chosen center or root branch approximates \(z\) to
\(Ka_1^{\mathrm{fit}}\), by \Cref{a08:local-test-bounds} and the
preceding center estimate.

There are \(K\) core/branch entries per label.  Delete light entries
again, using the same \(q_1\nu_\gamma(Q_1)\) budgets and a smaller
reciprocal-subpower threshold.  The total loss remains negligible,
and each retained entry has an earlier real projection of volume
at least \(1/K\).  Compose its center or root branch with the native
\(P_*\), using \(w=z/B\).  The native derivative and curvature
bounds imply an error at most \(Ka_1^{\mathrm{fit}}\) for \(X\)
at these incidences.  The resulting function is holomorphic and
algebraic of bounded relation degree; on the tested real projection
its size is at most \(K\), since \(|X|\le K_0\).
Apply \Cref{lem:algebraic-norm} on the comparison domains with
margin.  Its value and first derivatives on the used interiors
are bounded by \(K\).

The real cores lie strictly inside those interiors.  Boundary points
may be assigned to a neighboring core, or to a bounded overlapping
cover, without changing the \(K\) count.  Two points using the same
entry are thus compared within a derivative-controlled neighborhood;
different entries contribute separate graphs to the list.
All floors used to obtain these graph bounds precede conditioning
on a particular projected cell.  The bounds belong to the functions
and remain valid after further restrictions.


The core/branch norm bounds use the incidence witnesses retained after
the second deletion.  The earlier witnesses \(W_{\gamma,I,Q_1}\), kept
before it, remain available for counting whole compatible domains.
All deletions have negligible total mass by the original assignment
budgets.  We now fix \(\mu_{\mathrm{prep}}\); no further support
preparation will alter it.
\end{proof}

\subsection{The scalar projection: caps and support}

Scalar mesh projection requires two distinct estimates: a mass ceiling
for each scalar bin and a bound on the number of occupied bins.
A mass ceiling alone does not bound that number.  The graph
description just prepared will supply the support bound: we count the
earlier charts compatible with a projected cell, and then the scalar
bins met by their graphs.

Fix a world \(\gamma\), an \(A\)-block \(I\), and its rescaled time
\(\tau\in[0,1]\).  In this subsection the index \(A\) includes its
world and block whenever they are suppressed.  Put
\begin{equation}
 \label{a08:projected-coordinates}
 Z=e_Z\cdot y,\qquad
 \zeta=\frac{Z-Z_*}{q_A},
\end{equation}
where \(Z_*\) is the tangential component of its center at the block
midpoint.  Thus \(\zeta\) is affine with
\(d\zeta/d\tau=e_Z\cdot V\), while \(X=P_*\) is quadratic along each
trajectory.  Both have actually subpower ranges and coefficients on
the unit block.  Let \(T_A\) be this coordinate change on incidences,
retaining the trajectory data, and define
\begin{equation}
 \label{a08:A-law}
 \lambda_A
   =\frac{(T_A)_*(\mu_{\mathrm{prep},\gamma}|_{I,A})}
          {q_A\nu_\gamma(A)},\qquad
 R_A=\omega_\gamma q_A\nu_\gamma(A).
\end{equation}
The denominator is the fixed prior-times-block budget.  Thus
\(\lambda_A\) is a subprobability dominated by the rescaled
conditional trajectory prior \(\nu_\gamma(\,\cdot\,|A)\) times
\(d\tau\); it is not normalized by its current success mass.
This denominator remains unchanged under every later restriction.
Disjointness gives
the global identities and bound
\begin{equation}
 \label{a08:prepared-mass-budget}
 \|\mu_{\mathrm{prep}}\|
    =\sum_A R_A\|\lambda_A\|,\qquad \sum_A R_A\le1.
\end{equation}
The projected law means the pushforward of \(\lambda_A\) to
\((\tau,\zeta,X,\zeta')\).

\begin{lemma}[Projected cap and support bounds]
 \label{a08:projected-bounds}
 Put \(a=N^{-s_1}\).  The projection of the fixed-budget law
 \(\lambda_A\) in every used \(A\) has the following properties.
 \begin{enumerate}
  \item An \(X\)-bin of width \(p=N^{-r}\), \(0\le r\le s_1\), has
  density at most \(Kp^d\) per unit Lebesgue area in \((\tau,\zeta)\).
  At \(p=a\), adjoining a \(\zeta'\)-bin of width \(N^{-v}\),
  \(0<v\le u_1\), gains a factor \(KN^{-Sv}\), for some fixed
  \(0<u_1<u\).
  \item For a fixed sufficiently large \(M\), choose dyadic
  \(\sigma\asymp N^{-M}\).  Each \((\tau,\zeta)\)-mesh cell of width
  \(\sigma\) meets at most \(Ka^{-d}\) occupied \(X\)-bins of width
  \(a\).  In particular the total number of occupied
  \((\sigma,\sigma,a)\)-cells is at most
  \(Ka^{-d}\sigma^{-2}\).
 \end{enumerate}
\end{lemma}

\begin{proof}
\emph{Projected densities.}
At fixed original time, an interval of length \(b\) in \(\zeta\)
cuts area at most \(K_0J_Ab\) in the \(A\)-domain.  In unit spatial
coordinates of \(A\), it tests a linear form whose norm is
\(\|D_A^{\mathsf T}e_Z\|/q_A\ge1/K_0\), by
\Cref{a08:tangential-fullness}.  Thus this estimate is also a density
bound for the area pushforward.

At each exact original base, an \(X\)-bin of width \(p\) requires at
most \(K_0\) hidden bins of width \(p/B\), by the native derivative
lower bound and quadratic monotonicity.  Integrate the prepared pure
cap across the \(A\)-domain, change time by \(dt=q_A\,d\tau\), and
divide by \(q_A\nu_\gamma(A)\).  The resulting projected density is
bounded by
\[
 Kp^d\,\frac{J_A}{\nu_\gamma(A)}\le Kp^d,
\]
where the last \(K\) uses \Cref{a08:balanced-density} only on the
chosen diagonal.

For the angular refinement, the normal component of original \(V\)
is specified by \(C\) to accuracy \(K_0g_C\).  Once a \(\zeta'\)-bin
of width \(N^{-v}\) is specified, the full original \(V\) therefore
fits in \(K_0\) old bins of that width, for any fixed
\(0<v\le u_1<u\) and sufficiently large \(i\).
On the good sampled terminal entries, their joint numerator costs
at most \(KN^{-Sv}\) times the old pure terminal density.
Sum these denominators only within permitted good hidden-bin
ancestors at depth \(s_1\), then perform the same base integration.
This proves the angular cap.  No unmasked exceptional mass is
charged against the small angular factor.

\emph{Counting whole compatible domains.}
Fix \(A\), one projected cell, and a \(Q_1\)-time block meeting that
cell.  There are at most two such blocks, because \(q_A\sigma\ll q_1\).
We show that every compatible \(Q_1\)-domain in this world and block
lies in a common enlarged moving box of area \(K_0q_1^2g_1\).

At a compatible occurrence, its center has tangential distance at
most \(K_0q_1\) from the cell's \(Z\)-location.  Its normal distance
from the affine \(C\)-center is at most \(K_0q_1g_1\), using
\Cref{a08:C-Q1-comparison}, the true position bounds, and
\(g_C\ll q_1g_1\).  The normal center difference changes at rate at
most \(K_0g_1\) by \Cref{a08:C-Q1-comparison}.  The tangential center
difference changes at rate at most \(K_0\).  The same center bounds
therefore hold throughout the entire \(Q_1\)-block.

The frames obey the same constant angle comparison throughout
that block.  Projecting a major axis of length \(K_0q_1\) onto
\(n_C\) costs at most \(K_0q_1g_1\), and the minor axis already has
that size.  Thus the \emph{whole} domain of each compatible label,
not only its intersection with \(A\), lies in the claimed common
box.  The entire witness measure \(W_{\gamma,I,Q_1}\), whose mass
satisfies \Cref{a08:Q1-witness-floor}, can consequently be charged
there.

Those witnesses are disjoint in the original block and obey the
original zero-depth base cap.  If there are \(L_1\) compatible
labels, \Cref{a08:old-layer-masses,a08:Q1-witness-floor} give
\[
 \frac{L_1}{K}\,q_1 n_i(r_{1,i})q_1^2g_1
     \le Kq_1^3g_1.
\]
It follows that
\begin{equation}
 \label{a08:compatible-count}
 L_1\le \frac{K}{n_i(r_{1,i})}\le KN^{ds_1}.
\end{equation}
This is a separate bound for the fixed \(A\) and cell.  It does not
sum over all \(C\)-labels or use a cap conditional on \(A\).

\emph{Coordinate precision without an absolute-width loss.}
Fix one compatible \(Q_1\), with the normalized base coordinates
\(\widehat u\) of \Cref{a08:earlier-graphs}.  For two occurrences in the projected cell,
write \(\Delta t=t'-t\) and \(\Delta y=y'-y\), allowing the two
occurrences to be on different trajectories.  Their projected
coordinates give
\[
 |\Delta t|+|\Delta Z|\le K_0q_A\sigma.
\]
Since both are in the same \(C\)-domain,
\[
 n_C\cdot\Delta y
   =n_C\cdot y_C'\,\Delta t+O(K_0g_C).
\]
Decompose \(n_1=c_1n_C+c_2e_Z\), with \(|c_2|\le K_0g_1\), and
subtract the same \(Q_1\)-center at the two times.  The constant
velocity comparison in \Cref{a08:C-Q1-comparison} gives
\begin{align}
 \label{a08:normal-displacement}
 |n_1\cdot(\Delta y-y_1'\Delta t)|
 &\le K_0(g_1q_A\sigma+g_C),\\
 |\text{major component of }\Delta y-y_1'\Delta t|
 &\le K_0(q_A\sigma+g_C).
 \notag
\end{align}
Divide by the respective widths \(q_1g_1,q_1\), and include the
normalized time difference.  Thus
\begin{equation}
 \label{a08:relative-coordinate-error}
 |\Delta\widehat u|
 \le K_0\left(\frac{q_A\sigma}{q_1}
          +\frac{g_C}{q_1g_1}\right).
\end{equation}
The factor \(g_1\) in the first term of
\Cref{a08:normal-displacement} cancels before the narrow width is
inverted.

For clarity, write \(\Delta_i=\alpha_{2,i}-\alpha_{1,i}\).
The two terms in \Cref{a08:relative-coordinate-error}, including
subpower factors, are bounded on the diagonal by
\[
 N^{-M+ks_1+o(1)}
   +N^{-\Delta_i+2u+ks_1+o(1)}.
\]
Choose \(M>1+s_1+ks_1\), and then use \(\Delta_i\to\infty\).
Both terms are smaller than \(a=N^{-s_1}\) by a fixed power.
No fixed-power lower bound for \(g_1\) has been used.

Combine \Cref{a08:relative-coordinate-error} with the derivative and
approximation bounds in \Cref{a08:earlier-graphs}. Points of one
projected cell using one graph vary in \(X\) by at most \(Ka\).
Each graph therefore meets only \(K\)
\(X_a\)-bins.  Multiply by the \(K\) graphs per label and
\Cref{a08:compatible-count} to prove the cellwise support bound.
Finally the \((\tau,\zeta)\)-range has area at most \(K\), so the
total cell count is at most \(Ka^{-d}\sigma^{-2}\).
Both conclusions concern the same fixed law \(\lambda_A\).
They persist under every subsequent restriction.
\end{proof}

\subsection{Scalar candidates and an actual improvement}

We now construct a known atlas at the fixed positive depth
\(c'\) of \Cref{a08:gain}, with time length \(q_A\), actual
subpower lists and geometric slacks, and the coverage and tempered
complexity required by the finite-sample exclusion.

\begin{proof}[Proof of \Cref{prop:unbounded-improvement}]
We use the fixed prepared measure \(\mu_{\mathrm{prep}}\) from
\Cref{a08:earlier-graphs}, whose projections satisfy
\Cref{a08:projected-bounds}.  Residuals and covered pieces
below are submeasures of it; their \(A\)-normalizations always use
the denominator in \Cref{a08:A-law}.

\emph{A candidate on every substantial residual.}
Consider any subsequence and any residual
\(\mu_{\mathrm{res}}\le\mu_{\mathrm{prep}}\) of mass at least a
fixed positive fraction of \(\mathfrak p_{\mathrm{path}}\).
Set
\[
 \lambda_{A,\mathrm{res}}
   =\frac{(T_A)_*(\mu_{\mathrm{res},\gamma}|_{I,A})}
          {q_A\nu_\gamma(A)}\le\lambda_A.
\]
The identity in \Cref{a08:prepared-mass-budget}, applied to this
residual, shows that some \(A\) satisfies
\(\|\lambda_{A,\mathrm{res}}\|\ge1/K\).
Use the fixed conditional prior \(\nu_\gamma(\,\cdot\,|A)\) and
active law \(\lambda_{A,\mathrm{res}}\) in
\Cref{thm:scalar-mesh-projection}.  The two mesh hypotheses, with
\(c=1\), follow from \Cref{a08:projected-bounds}:
\[
 \sigma^2
 \#\{\text{occupied }(\sigma,\sigma,a)\text{-cells}\}
       \le Ka^{-d},
\]
together with the pure caps and a positive terminal angular gain.
The upper caps and occupied-cell bound persist under this
restriction, and the displayed mass floor supplies the required
density relative to the fixed prior.

On a further subsequence the theorem gives a quadratic
\(F(\tau,\zeta)\) fitting \(X\) to \(Ka\) on assigned trajectories
throughout the unit block, with residual incidence mass in original
units at least
\begin{equation}
 \label{a08:candidate-floor}
 \frac{q_A\nu_\gamma(A)a^d}{K}.
\end{equation}
Indeed the angular conclusion makes its moving interval scale
\(\beta\sim_{\log,N}1\).  The theorem's adapted norm bound then
bounds the ordinary coefficients of \(F\) by \(K\): the interval
center motion and all occupied coordinate ranges are bounded by
\(K\).

\emph{Greedy coverage in original mass.}
In every \(A\), greedily insert quadratic tests whose coefficients
are at most \(N^\delta\), whose whole-block fit is at most
\(aN^\delta\), and whose residual masked incidence mass is at least
\begin{equation}
 \label{a08:greedy-floor}
 q_A\nu_\gamma(A)a^dN^{-\delta}.
\end{equation}
Let \(\mu_{\mathrm c}\le\mu_{\mathrm{prep}}\) denote the incidences
covered by these assignments, measured with their original world
weights; at any stage the residual is
\(\mu_{\mathrm{prep}}-\mu_{\mathrm c}\).
On insertion assign \emph{all} still-unassigned \(A\)-trajectories
meeting the geometric fit.  These assignments are invariant over
the block.  Each insertion removes at least
\(a^dN^{-\delta}\nu_\gamma(A)\) trajectory mass, since its
incidence duration is at most \(q_A\).  Therefore there are at most
\(N^\delta a^{-d}\) inserted tests in each \(A\).

Maximality forces substantial coverage on a subsequence.  If the
covered fraction \(\|\mu_{\mathrm c}\|/\mathfrak p_{\mathrm{path}}\)
tended to zero, the uncovered incidences would be a
substantial residual of the prepared path.  The preceding
application of \Cref{thm:scalar-mesh-projection} would produce
another candidate on a further subsequence.  Since \(K\le N^\delta\)
eventually, its coefficient, fit and mass bounds would make it
eligible in \Cref{a08:greedy-floor}, contradicting termination.
Thus for some fixed \(\kappa>0\), along a subsequence,
\[
 \|\mu_{\mathrm c}\|\ge\kappa\mathfrak p_{\mathrm{path}}.
\]
An individual candidate's reciprocal-subpower floor is used to
bound the insertion count, not as the eventual coverage.

\emph{Posterior comparison.}
Keep \(\mu_{\mathrm c}\) unnormalized, including its analytical
masks.  Let \(O(\xi)=(p_1(\xi),A(\xi))\) be the observation of an
incidence \(\xi\), with its world and block included, and put
\(m_{\mathrm c}=O_*\mu_{\mathrm c}\).
For \(m_{\mathrm c}\)-almost every observation \(o\), let
\(\pi_o\) be the conditional probability on incidences.  Define the
posterior-pair measure by
\begin{equation}
 \label{a08:posterior-pair-law}
 \Pi=\int (\pi_o\otimes\pi_o)\,dm_{\mathrm c}(o).
\end{equation}
The base measure in this integral is the unnormalized observation
pushforward.  Consequently \(\|\Pi\|=\|\mu_{\mathrm c}\|\), and
both incidence marginals of \(\Pi\) equal \(\mu_{\mathrm c}\).
The two assigned polynomials \(F,G\) are evaluated at the same
\((\tau,\zeta)\).  At the common original base their hidden
coordinates lie in one terminal bin; the native derivative and
curvature bounds give
\[
 |X_F-X_G|\le K_0N_0^{-1}.
\]
Since \(s_1<1\), the two fits consequently imply
\begin{equation}
 \label{a08:posterior-match}
 |F(\tau,\zeta)-G(\tau,\zeta)|
        \le aN^{O(\delta)}.
\end{equation}

Use the maximum norm on the six quadratic coefficients in the
original \((\tau,\zeta)\)-coordinates, and define
\[
 \mathcal B
   =\{(\xi_1,\xi_2):\|\operatorname{coeff}F-
                              \operatorname{coeff}G\|_\infty
                              >a^{2\theta}\},\qquad
 \Pi_{\mathrm{bad}}=\1_{\mathcal B}\Pi.
\]
We claim that \(\|\Pi_{\mathrm{bad}}\|=o(\|\Pi\|)\).
Otherwise it has a fixed positive relative mass on a subsequence.
The budgets in \Cref{a08:A-budget} then give a world and an
\(A\)-block for which the within-world restriction
\(\Pi_{\mathrm{bad},\gamma,A}
 =\omega_\gamma^{-1}\Pi_{\mathrm{bad}}|_{\gamma,I,A}\) has
\[
 b_A:=\|\Pi_{\mathrm{bad},\gamma,A}\|
       \ge\frac{q_A\nu_\gamma(A)}{K}.
\]
Normalize only these selected pairs, after the coordinate change:
\begin{equation}
 \label{a08:bad-pair-normalization}
 \widehat\Pi_A
   =\frac{(T_A\times T_A)_*\Pi_{\mathrm{bad},\gamma,A}}{b_A},
 \qquad
 (\widehat\Pi_A)_j
   \le\frac{q_A\nu_\gamma(A)}{b_A}\lambda_A
   \le K\lambda_A\quad(j=1,2).
\end{equation}
Here \((\widehat\Pi_A)_j\) is the \(j\)-th incidence marginal.
The domination follows because each marginal before normalization
is bounded by \(\mu_{\mathrm c,\gamma}|_{I,A}\), hence by
\(\mu_{\mathrm{prep},\gamma}|_{I,A}\).

For this probability law each individual polynomial index has
marginal base density at most
\begin{equation}
 \label{a08:index-base-ceiling}
 KN^{O(\delta)}a^d
 \quad\text{in }(\tau,\zeta).
\end{equation}
Indeed, at a fixed base the fit of one assigned polynomial permits
only \(N^{O(\delta)}\) \(X_a\)-bins.  Their caps under
\(\lambda_A\), followed by \Cref{a08:bad-pair-normalization}, give
this ceiling.  Each color has a list of at most
\(N^\delta a^{-d}\) indices.  If \(\rho_{j,f}\) is uniform on that
list, then \(a^d\le N^\delta\rho_{j,f}\); hence the index-base
marginal is bounded by
\(KN^{O(\delta)}\rho_{j,f}\,d\tau\,d\zeta\), as required for
local graph matching.

Enlarge and rescale the \(\zeta\)-range by a factor between one and
\(K\) into a fixed bounded comparison domain.  The polynomial norm
upper bounds become \(KN^\delta\), while a coefficient discrepancy
larger than \(a^{2\theta}\) gives a norm gap bounded below by a
constant times that quantity, allowing subpower comparison factors.
Apply \Cref{a02:local-graph-matching} to these two-variable
quadratic graphs, which are holomorphic on any fixed larger
comparison domain, with common reference zero, norm scale
\(M=1\) and matching scale \(\Delta=a\); thus \(R=a^{-1}\).
Use the tolerance \(\eps\) fixed before \Cref{a08:gain}.
For every fixed loss exponent \(C\le C_{\mathrm{match}}\),
\[
 N^{C\delta}=R^{C\delta/s_1},\qquad
 C\delta/s_1<\eps/2.
\]
This accounts for the list count, polynomial norm, value discrepancy
and index-base density.  The range rescaling and the bad-pair
normalization contribute only factors \(K=R^{o(1)}\), so all these
upper bounds fit the tolerance \(\eps\) at late samples.
The norm gap is at least \(R^{-2\theta}/K\), which exceeds
\(R^{-\eps}\) because \(2\theta<\eps\).
Finally \(\widehat\Pi_A\) has mass one and \(d<d_0<1\).
Thus all hypotheses of \Cref{a02:local-graph-matching} hold,
contradicting its conclusion.  This proves the claim about
\(\Pi_{\mathrm{bad}}\).

Bin the six coefficients of a quadratic in \((\tau,\zeta)\), in
these original \(A\)-coordinates before the range rescaling, to
width \(a^\theta\).  The successful posterior pairs lie in bounded
neighboring bin lists, because \(a^{2\theta}\ll a^\theta\).
For an observation \(o=(p_1,A)\), let \(\beta_o\) be the
coefficient-bin pushforward of the incidence posterior \(\pi_o\),
and let \(L(b)\) be the bounded neighboring list around a bin \(b\).
For independent bins \(B_1,B_2\) with law \(\beta_o\), the matching
event is \(B_2\in L(B_1)\), and
\[
 \int \beta_o(L(b))\,d\beta_o(b)
   =\Pp\{B_2\in L(B_1)\mid o\}.
\]
Choose a maximizing center \(b\) for each \(o\), and integrate
against \(m_{\mathrm c}\).
This gives deterministic bounded lists retaining a fixed positive
fraction of the covered mass.  Composing with the actual
\(A\)-lists at \(p_1\) multiplies list cardinality only by an
actual subpower factor.

\emph{The new test and its actual bounds.}
Within each \(A\), merge all assigned tests with the same coefficient
bin. For a bin \(b\), let
\(\widetilde F_{A,b}\) be the quadratic with its binned coefficients.
The final chart label is \((A,b)\).
In original coordinates use
\begin{equation}
 \label{a08:final-test}
 P_{\mathrm{new}}(t,y,w)
   =P_*(t,y,w)
    -\widetilde F_{A,b}\left(
       \frac{t-t_A}{q_A},
       \frac{e_Z\cdot y-Z_*}{q_A}\right),
\end{equation}
where \(t_A\) is the block's left endpoint.
The error on every assigned trajectory throughout the block is
at most
\[
 aN^\delta+Ka^\theta.
\]
Our fixed choices give \(\delta<\theta s_1/4\),
\(\theta<1\) and \(c'=\theta s_1/2\).  Thus both exponents
\(s_1-\delta\) and \(\theta s_1\) are strictly larger than \(c'\),
so at late diagonal samples the error is at most \(N_0^{-c'}\).
The positive power buffer absorbs the merely diagonal subpower
factor \(K\).

The polynomial subtracted in \Cref{a08:final-test} does not involve
\(w\).  Thus
\[
 (P_{\mathrm{new}})_w=(P_*)_w,\qquad
 (P_{\mathrm{new}})_{ww}=(P_*)_{ww}.
\]
The original native bounds give the actual full-block derivative
upper bound and, on retained original path incidences, the actual
lower bound.  The native curvature bound \(K_0B^2\) is stronger
than the required \(K_0B^2N_0^{c'}\).
Keep the actual frame and time block of \(A\).  All geometric
slacks and derivative requirements therefore cost only fixed
powers of actual path factors.  The time exponent still tends
to zero.

\emph{Assignments and lists on the original law.}
Store the ordered quadratic tests for each \(A\), assign trajectories
by the first geometric whole-block fit, and merge tests with the same
coefficient bin. Each whole-block fit is a fixed-degree semialgebraic
condition on the trajectory coefficients, obtained by eliminating its
single time quantifier. The at most \(N^\delta a^{-d}\) tests in each
\(A\), their priorities and merger table therefore have polynomial
complexity. The coordinate changes, coefficients and inverse lengths
also have polynomial bounds within each fixed \(i\).

Attach this geometric label map to the original tempered path
\(\mu_{\mathrm{ref}}\), using a failure label on unassigned pieces.
It agrees with the posterior construction on \(\mu_{\mathrm c}\).
The latter measure has the bounded coefficient-bin lists just proved,
composed with the actual \(A\)-lists at \(p_1\), and
\[
 \mu_{\mathrm c}\le\mu_{\mathrm{prep}}\le\mu_{\mathrm{ref}}.
\]
Thus \Cref{a03:actual-lists} applies with original law
\(\mu_{\mathrm{ref}}\), analytical submeasure \(\mu_{\mathrm c}\),
output label \((A,b)\), and original observation
\(p_1\). The failure label is omitted from every list. In particular,
the lemma uses the unnormalized success mass already obtained; no
comparison of normalized posteriors is needed.

All predicates entering this application are the original path
predicates and the finitely many geometric fit tests. The analytical
cap masks and residual searches are absent. The polynomial bounds
for these geometric data give a fixed mesh exponent and list table
of the complexity allowed by the prechosen finite-sample exclusion. The
table loses only \(O(N_0^{-1})\) mass, and its list sizes remain fixed
powers of actual path factors and logarithms.

A fixed fraction of \(\kappa\mathfrak p_{\mathrm{path}}\) therefore
survives. Since \(\mathfrak p_{\mathrm{path}}\ge N_0^{-o(1)}\),
the flattening error is negligible and the resulting coverage exceeds
\Cref{a08:required-output-mass} for all sufficiently late samples.
Together with the geometric and derivative bounds above, this is
precisely the excluded known atlas at depth \(c'\), with time
\(q_A\). It contradicts \Cref{a08:forbidden-horizon} and proves
\Cref{prop:unbounded-improvement}.
\end{proof}

\subsection{Completion of the proof}
\label{proof:main}

\begin{proof}[Proof of \Cref{thm:main}]
Let \(K\subset\R^4\) contain a unit line segment in every direction.
Suppose that \(\dim_{\mathrm H}K<4\), and choose \(\sigma>0\) with
\(\dim_{\mathrm H}K<4-\sigma\).
By \Cref{thm:model-reduction}, the Hausdorff covers of this arbitrary
set produce an exact problem of the second model in \(C(d_0)\),
for some \(d_0<1\) and fixed positive angular parameters, with
\(H>0\).  That reduction uses outer slope measure and finite
selections of witnesses, so it applies to nonmeasurable \(K\) and
nonmeasurable witnessing line families.

Optimize cap dimension and horizon as in \Cref{prop:critical-path}.
The resulting critical dimension is less than one and the
maximizing time rate satisfies \(k>0\).
The auxiliary first model has already been excluded by
\Cref{a05:scalar-zero-horizon}; this supplies the scalar results
used in the subsequent second-model branches.

For the second model, consider the global optimized narrowness
\(\ell\).  If \(\ell=0\), \Cref{thm:isotropic-improvement} gives a
forbidden improvement.  For \(0<\ell<\infty\), the projection and
comparison arguments in
\Cref{prop:wide-interval,prop:horizontal-model} either give such an
improvement immediately or reduce to the horizontal case with
narrowness rate equal to \(k\).
At \(2k=1\), \Cref{prop:critical-rate} excludes that case.
Above the critical rate, \Cref{prop:offcritical-wide} supplies the
wide intervals required by \Cref{prop:wide-interval}.
Below the critical rate, \Cref{prop:bootstrap} first reaches the
limiting parameter scale, and \Cref{prop:offcritical-wide} then
supplies those intervals.
In each instance the original-unit construction of
\Cref{prop:candidate-upgrade} yields the actual atlas forbidden
by the critical-path optimization.
Finally, the global case \(\ell=\infty\) is excluded by
\Cref{prop:unbounded-improvement}.

Every possible optimized narrowness has therefore been excluded,
contradicting the positive-horizon problem supplied by
\Cref{thm:model-reduction}.  It follows that
\(\dim_{\mathrm H}K\ge4\).  The ambient upper bound is
\(\dim_{\mathrm H}K\le4\), so \(\dim_{\mathrm H}K=4\).
The argument imposes no compactness, measurability or stickiness
condition on \(K\) or on its witnessing line family.
\end{proof}
